\documentclass[10pt,a4paper]{article}
\usepackage[utf8]{inputenc}
\usepackage{amsmath}
\usepackage{amsfonts}
\usepackage{comment}
\usepackage{amssymb}
\usepackage{graphicx}
\usepackage[left=2cm,right=2cm,top=2cm,bottom=2cm]{geometry}
\usepackage{listings}
\usepackage{color}

\usepackage{pgfplots}
\pgfplotsset{compat=1.15}
\usepackage{mathrsfs}
\usetikzlibrary{arrows}
\pagestyle{empty}


\definecolor{munsell}{rgb}{0.0, 0.5, 0.69}
\definecolor{dkgreen}{rgb}{0,0.6,0}
\definecolor{gray}{rgb}{0.5,0.5,0.5}
\definecolor{mauve}{rgb}{0.58,0,0.82}
\definecolor{minas}{RGB}{0.244, 0.172, 0.36}
\definecolor{PUJ}{RGB}{44, 86, 151}
\definecolor{PUJ2}{RGB}{43, 93, 156}
\definecolor{PUJ3}{RGB}{20, 52, 107}
\definecolor{cyandk}{rgb}{0.0, 0.72, 0.92}

\author{Iván Andrés Trujillo }

%% code information
\lstset{frame=tb,
  language=Python,
  aboveskip=3mm,
  belowskip=3mm,
  showstringspaces=false,
  columns=flexible,
  basicstyle={\small\ttfamily},
  numbers=none,
  classoffset=1,
  morekeywords={True,False}, keywordstyle=\color{munsell}, 
  classoffset=0, 
  keywordstyle=\color{blue},  
  commentstyle=\color{dkgreen},
  stringstyle=\color{PUJ3},
  numberstyle=\tiny\color{gray},
  breaklines=true,
  breakatwhitespace=true,
  tabsize=4,
}

\begin{document}

\newcommand{\HRule}{\rule{\linewidth}{0.5mm}}

%----------------------------------------------------------------------------------------
%	LOGO SECTION
%----------------------------------------------------------------------------------------

\includegraphics[width = 4cm]{pujshield.eps}\\[0.5cm] 

\begin{center} 
\textbf{\LARGE Elementary mathematics for machine learning }\\[0.2cm]
\emph{\LARGE Notes of class}\\[0.3cm] 
\emph{Iván Andrés Trujillo Abella \\
Cesar agusto del corral } \\
\textsc{\Large 
}\\[0.2cm] 
\textsc{\large Facultad de ingeniería}\\[0.5cm] 
\HRule \\[0.4cm]
\end{center}
\vspace{1cm}



\subsection{Jacobian Matrix}

linear maps or linear transformations:



a alternative coordinate system is related with jacobian.

Rememeber that if you alter the base coordinate vectors



note 

jacobian determinant: how much areas are scale near to one point.




here is important the following:
we defined that is easy find a analytical transformation only indicating how many times 

$\alpha_{j} \hat{j} + \alpha_{i} \hat{i}$ this allow us determined the new vector in linear transform, and the linear transform must be 
have the following properties:

\begin{equation}
\begin{align*}
T(k\vec{v}) &= k \\
T(\vec{v} + \vec{u} ) &= T(\vec{v}) + T(\vec{u})
\end{align*}
\end{equation}


\begin{comment}
warped:
 $e_{1} = (0,1) , e_{2}=(1,0)$ 
vanish
rod:
gaugue:

\end{comment}




WHat is base?
what is a subspace?



\subsection{vector space}



\subsection{Propierties of vector spaces}

Conmutativity $\vec{u} + \vec{v} = \vec{v} + \vec{u}$, null vector, identity element and inverse ( $\vec{v} + -\vec{v} )  = 0$

\subsection{Properties of scalars}

multiple a vector $\vec{u}$ of $n-components$ by a scalar $k$  where $k \in R $ then we get $k\vec{u} = (ku_{1}, ku_{2},...,ku_{n})$.




\section{ideas or insights}
adding two scaled vectors.
the unitary vectors $\vec{i}$ and $\vec{j}$ the unitary vectors are the basis vectors of the coordinate system.


\section{linear combination}

it is called linear combination, to the following 
\begin{equation}
k_{1}\vec{i} + k_{2} \vec{j}
\end{equation}
note that $k_{1},k_{2} \in R$ also take in account that if we  fix a vector and vary the another we can get a straight line.

if you take in mind that $\vec{i},\vec{j}$ it is measure of length then we can get the point $(x,y)$ as the sum of the vectors by the scalars $x,y$.




\subsection{Independence basis and dimension}


\subsection{Positive define matrices and co-variance matrix  }

think in the following problem, in OLS you must 


think in the following problem  in a quadratic form we must sure that the exist a minimun point; 



\begin{equation}
ax^{2}+ by^{2}+cxy
\end{equation}

is this function positive ?
note that the squares are also be positive however $cxy$ could be negative the question is that $ax^{2}+by^{2} - cxy >0$ always.


we said that a positive definite (or semi-definite) matrix $\textbf{A}$ if 

\begin{equation}
\textbf{x}^{T} \textbf{Ax} \geq 0, \forall  \textbf{x} \in R^{n}
\end{equation}
if the inequality is strictly then is definite otherwise is semi.






\section{span}
is the set of linear combinatios of two vectors i


\section{plane}
two vectors span a plane.
three vectors span a $R^{3}$.



\section{Basis for a subspace}
independent vectors that span the subspace.


\section{Dimension}
number of basis vectors for the subspace.




\definecolor{xdxdff}{rgb}{0.49019607843137253,0.49019607843137253,1}
\begin{tikzpicture}[line cap=round,line join=round,>=triangle 45, scale=1]
\begin{axis}[
axis lines=middle,
yticklabels={},
xticklabels={},
x=4cm,y=4cm,
xmin=-0.3,
xmax=3.5,
ymin=-0.3,
ymax=2.4]
\clip(-0.4917288227170077,-0.5202856575366671) rectangle (4.12143238498744,2.444455380399042);
\draw [->,line width=1pt] (0,0) -- (1,0);
\draw [->,line width=1pt] (0,0) -- (0,1);
\draw (0.44418739088795206,-0.025155789694043584) node[anchor=north west] {$\mathbf{\vec{i}}$};
\draw (-0.21397304319553578,0.6420618980694918) node[anchor=north west] {$\mathbf{\vec{j}}$};
\draw [->,line width=1pt] (0,0) -- (3,0);
\draw (1.8903288951356156,-0.007041282333947601) node[anchor=north west] {$\mathbf{x\vec{i}}$};
\draw [->,line width=1pt] (0,0) -- (0,2);
\draw (-0.1868012821553918,1.6353407183147548) node[anchor=north west] {$\mathbf{y\vec{j}}$};
\draw [->,line width=1pt] (0,0) -- (3,2);
\draw [line width=1pt,dotted] (0,2)-- (3,2);
\draw [->,line width=1pt] (0,0) -- (3,0);
\draw (2.9319130683411356,2.1757568545576182) node[anchor=north west] {$(x,y)$};
\begin{scriptsize}
\draw [fill=xdxdff] (3,2) circle (2pt);
\end{scriptsize}
\end{axis}
\end{tikzpicture}


















\section*{Complex numbers}
a complex number is composed of a real part and imaginary, $z = r+i$ , the imaginary part solve $i^{2}= -1$.



\section{Eigen values and eigen vectors}

take in mind the following operations:
\begin{itemize}
\item translation
\item inversion
\item projections
\item stretch 
\item rotation
\end{itemize}


\begin{equation}
\textbf{A}\vec{x}  = \lambda \vec{x}
\end{equation}

Note that the map of $A$ over $\vec{x}$ 




\subsection{Translation}
Only consist in pass from the coordinates $(x,y)$ to 
$(x+\delta x, y)$ 

\[
\begin{pmatrix}
x + \delta x \\
y 
\end{pmatrix}
=
\begin{pmatrix}
w & v \\
z & j
\end{pmatrix}
\begin{pmatrix}
x\\
y
\end{pmatrix}
\]
therefore; 


\[
\begin{pmatrix}
x + \delta x \\
y 
\end{pmatrix}
=
\begin{pmatrix}
1 + \delta & 0 \\
0 & 1
\end{pmatrix}
\begin{pmatrix}
x\\
y
\end{pmatrix}
\]

Note that the translation is over the $x-axe$ over the $y-axe$ is only $(x,y)$ to $(x,y + \delta y ).$
\subsection{Rotation Matrix}

Rotation matrices, rotation is based on the idea of projections.

remember that a vector in a coordinate system (A)  could be written as $\vec{x}^{A}$ also remember that in a analytic way.

in 2D for instance $\vec{x} = a\vec{i} + b\vec{j}$ where $\vec{i},\vec{j}$ are unitary vectors, in 3D is equal
\begin{equation}
\vec{x}^{A} = \begin{pmatrix}
k_{1}\vec{i}^{A} \\
k_{2} \vec{j}^{A} \\
k_{3} \vec{z}^{A}
\end{pmatrix}
\end{equation}

how calculate the rotation from a coordinate A to coordinate B.




\definecolor{ududff}{rgb}{0.30196078431372547,0.30196078431372547,1}
\definecolor{xdxdff}{rgb}{0.49019607843137253,0.49019607843137253,1}
\definecolor{uuuuuu}{rgb}{0.26666666666666666,0.26666666666666666,0.26666666666666666}
\begin{tikzpicture}[line cap=round,line join=round,>=triangle 45,x=1cm,y=1cm]
\begin{axis}[
x=1cm,y=1cm,
axis lines=middle,
xmin=-2.4496946307638305,
xmax=8.761019121625118,
ymin=-0.9998544303848308,
ymax=6.204936737708068,
xtick={},
ytick={},]
\clip(-2.4496946307638305,-0.9998544303848308) rectangle (8.761019121625118,6.204936737708068);
\draw [->,line width=2pt] (0,0) -- (0,5.181952225161589);
\draw [->,line width=2pt] (0,0) -- (5,0);
\draw [->,line width=2pt] (0,0) -- (4,3);
\draw [->,line width=2pt] (0,0) -- (-2,5);
\draw (4.637706864529183,-0.09742132073572028) node[anchor=north west] {$\vec{i}^{A}$};
\draw (3.639894645730166,2.8079730810614163) node[anchor=north west] {$\vec{i}^{B}$};
\draw (0.12554131628363074,5.023703155159232) node[anchor=north west] {$\vec{j}^{A}$};
\draw (-1.7526934485145176,4.942997755109312) node[anchor=north west] {$\vec{j}^{B}$};
\begin{scriptsize}
\draw [fill=uuuuuu] (0,0) circle (2pt);
\draw [fill=xdxdff] (0,5.181952225161589) circle (2.5pt);
\draw [fill=ududff] (5,0) circle (2.5pt);
\draw [fill=ududff] (4,3) circle (2.5pt);
\draw [fill=ududff] (-2,5) circle (2.5pt);
\end{scriptsize}
\end{axis}
\end{tikzpicture}


Note that for this case, we pass from coordinates  A to coordinates B, we can denote this in $R^{A}_{B}$.

Note that we have project $\vec{i}^{B}$ over the vectors $\vec{i}^{A}$, and $\vec{j}^{A}$.

Note here that the coordinates $(x,y)$ now will be $(\hat{x},\hat{y})$ think in any vector generated by the translation$\vec{u} = k_{1}\vec{j}^{B} + k_{2}\vec{i}^{B}$, the question  is $\vec{u} = k_{1}\vec{j}^{A} + k_{2}\vec{i}^{A}$ what is the vector in terms of the A coordinate system.
\\

Assume that the vectors are rotate in $\theta$ degrees \footnote{
Remember that $\sin(\theta) = \frac{opposite}{Hypothenuse}$ and $\cos(\theta) = \frac{adjacent}{Hypothenuse}$.
}.

thus we can translate the projections over the system A, as follow:


\begin{equation}
\vec{i}^{B} = proyection(\vec{i}^{B})_{x} + proyection(\vec{j}^{B})_{y}
\end{equation}

Notice taht 
Now:

\begin{equation}
cos(\theta)= \frac{proyection(\vec{i}^{B})_{x}}{\vec{i}^{B}} = \cos(\theta)
\end{equation}

given that the vector $\vec{i}^{B}$ it is unitary. 


now 

\begin{equation}
proyection(\vec{i}^{B})_{y} = \sin (\theta)
\end{equation}

thus we can get:
\begin{equation}
\vec{i}^{B} = \cos (\theta) + \sin(\theta).
\end{equation}

the same is applied to the 
$\vec{j}^{B}$ getting:

\begin{equation}
\vec{j}^{B} = 
\end{equation}


\subsection{Getting eigen values}

think in mind the following, the span is the entire line that lay over the value.












\begin{comment}

countless:
crowded: atestado.
shaky:
knocked off:
sneakier:
glimpse: vistazo



\end{comment}




things to think, why $\vec{i}$ and $\vec{j}$ are the respectively columns and when you multiply, you multiply in  inner product.


rotations and projections how think in this operations and how are related?




\section{Linear approximation}

If we have $ y = f(x)$ the linear approximation is constructed with tangent line, with in $z = f(x,y)$ then we can think in a tangent plane.


\subsection{Quadratic approximation}
Is a better approximation (another curve) that a tangent plane.



\section{Hessian Matrix}
Second derivates of a function: 






\end{document}
